\chapter{Practical application guide}
\label{chap:guia}

\section{A reproducible first walkthrough}
QSB provides a macOS application and a command-line tool called \texttt{qsb}.
Both share the models, validators, and solvers in \texttt{QSBCore}. The current
interface uses English labels; this manual retains those labels so they can be
located on screen.

To begin, open \path{manual/examples/lp_sample.json} using
\textbf{Model > Open Model...}. The model maximizes $50X_1+60X_2$ subject to
$2X_1+3X_2\leq180$, $3X_1+2X_2\leq150$, and nonnegative variables.
\begin{enumerate}
 \item In \textbf{Definition}, check names, coefficients, relations, bounds,
 and variable types. The variables here are continuous.
 \item Click \textbf{Validate} and review \textbf{Validation}. A valid model
 is a consistent input; it is not yet a solution.
 \item In \textbf{Run}, select \textbf{QSB Native} and click the \textbf{Run}
 button; you can also use \textbf{Solve} in the toolbar or
 \textbf{Solve > Run Current Model} in the menu.
 \item In \textbf{Summary}, check $X_1=18$, $X_2=48$, and objective 3780.
 Both constraints hold with equality. Also review the solution status and
 metadata before interpreting the number as an optimum.
 \item Export the model and solution separately. Reopen the exported model
 to preserve a reproducible input.
\end{enumerate}
The \textbf{Samples} menu provides built-in examples, but these may differ from
the manual's files. To reproduce a cited result, use the file specified in its
chapter.

\section{Creating, importing, and editing}
\textbf{New Model...} opens the family selector and creates an initial model.
\textbf{Open Model...} accepts normalized JSON and recognized legacy formats.
The family determines which controls appear. Native editors let you work with
names, tables, and typed fields; other variants use a structured form and JSON
as an interchange surface.

The \textbf{Definition} panel defines the problem; \textbf{Overview} summarizes
its family and state. \textbf{JSON} lets you inspect its representation. When
\textbf{Apply JSON to ... Editor} is available, explicitly apply the edited text
to replace the native draft. Writing JSON does not mean the changes have
already been incorporated into the table. Review the diagnostics if applying the JSON fails and preserve the text you need to repair.

Some tables offer \textbf{Copy Table} and \textbf{Paste Table} to exchange ranges
with spreadsheets. Check headers, column order, separators, units, and names
after pasting. Malformed numeric fields are not silently converted to zeros.
In editors that flag incomplete imported dimensions, review the cells and
shown values before accepting \textbf{Use Displayed Dimensions}.

\section{Validating, running, and interpreting states}
The application distinguishes empty, editing, validating, invalid, and valid
model states. Execution has separate states: \textbf{Not run},
\textbf{Validating}, \textbf{Solving}, \textbf{Solved}, and \textbf{Failed}.
Read a Solved status alongside the mathematical condition reported by the
family: a method may return an approximation, local solution, or heuristic
policy rather than a global optimum.

\textbf{Validation} displays model errors and warnings. Correct errors before
solving; a warning requires judgment even when it allows you to proceed.
\textbf{Run Details} lets you review the backend and execution information.
When a long run displays progress, wait for completion before repeating Run.
There is no universal cancellation control; some editors' asynchronous paths
prevent an old response from replacing a new draft. This does not cancel the
background computation.

\begin{description}
 \item[QSB Native] Deterministic educational methods. Exactness and limits
 depend on the family and are explained in the corresponding chapters.
 \item[Validate only] Checks the input and produces diagnostics; it does not
 compute decisions or optimal values.
 \item[External solver] Integrations available only for certain families.
 Availability depends on configuration and the installed solver. The presence
 of this option does not imply universal external support.
\end{description}
For LP/ILP and QP, consult the HiGHS conditions in the optimization chapter.
Other families may report that the external backend is unavailable. Switch to
QSB Native or Validate only according to the work you want to do.

\section{Saving models and results}
The toolbar's \textbf{Export} button prepares JSON from the current draft. If a
native field cannot be converted into the model, QSB displays validation rather
than exporting old data. The \textbf{Model > Export Model JSON...} menu opens
export of the model text; to ensure recent fields are synchronized, first use
Validate or the toolbar's Export button. \textbf{Export Solution JSON...} saves
the existing result and is enabled when solution content is present.

These are distinct documents: the model describes the problem to solve, and the
solution records what a run returned. Output formats vary by family; some
include the model and others do not. Always preserve both, together with the
command, backend, software version, and study assumptions. Do not use a solution
document as a model input.

\section{Keyboard controls and alternatives to charts}
The application defines the following shortcuts:
\begin{center}
\begin{tabular}{ll}
\toprule
Action & Keys\\
\midrule
New Model & Cmd+N\\
Open Model & Cmd+O\\
Export Model JSON & Cmd+S\\
Export Solution JSON & Cmd+Shift+S\\
Validate Model & Cmd+Shift+V\\
Run Current Model & Cmd+R\\
Solve ILP & Cmd+I\\
\bottomrule
\end{tabular}
\end{center}
Buttons are enabled according to family and state. Solve ILP and Solve LP
Relaxation are specific to linear/integer programming; the relaxation does not
enforce integrality. Some editors use Cmd+Shift+N to add rows. Use Tab to move
through native controls and the menus to locate actions; focus behavior depends
on the editor.

Tables and JSON let you inspect data without relying on chart colors. Native
results include accessible labels and values in several families, but this
alone does not establish a complete VoiceOver review. Preserve the numerical
tables for reproducible analysis. The macOS window has a minimum size of 980 by
620 points; use scrollable areas and close the Inspector when you need more width.

\section{The terminal and normalized JSON}
From the repository root, SwiftPM builds and runs the tool:
\begin{lstlisting}
swift run qsb
swift run qsb validate manual/examples/lp_sample.json
swift run qsb solve-json manual/examples/lp_sample.json --backend native
\end{lstlisting}
Without arguments, \texttt{qsb} displays the command catalog. After
\texttt{swift build}, you can run \texttt{.build/debug/qsb} directly. To save
output and separate errors:
\begin{lstlisting}
.build/debug/qsb solve-json manual/examples/lp_sample.json \
  --backend native > /tmp/qsb-lp-solution.json 2> /tmp/qsb-lp-errors.txt
\end{lstlisting}
JSON uses decimal points, double-quoted strings, and \texttt{null} for optional
fields that allow it. It does not include comments. Each family has its own
schema: LP uses a direct object, while many other families use a wrapper with
\texttt{kind} and \texttt{model}. Copy the structure of the family's example,
keeping names, vectors, and matrices aligned.

\texttt{solve-json} detects the family of a normalized input.
\texttt{validate} accepts normalized or legacy inputs. \texttt{solve} targets
legacy models; do not interchange commands indiscriminately. The chapters
provide family-specific commands, such as \texttt{solve-dp-json} or
\texttt{solve-decision-json}. \texttt{validate-json} is specific to LP/ILP,
rather than a universal validator.

To import your own legacy input, the general procedure is:
\begin{lstlisting}
qsb import-legacy-json /ruta/al/modelo > /tmp/qsb-model.json
qsb validate /tmp/qsb-model.json
qsb solve-json /tmp/qsb-model.json --backend native
\end{lstlisting}
Here \texttt{qsb} assumes that the executable is on PATH; you can replace it
with \texttt{swift run qsb}. The original files must be locally accessible and
used in accordance with their permissions. The manual's examples are
independent of proprietary material.

\section{Troubleshooting}
\begin{description}
 \item[Unreadable JSON] Check syntax, key names, and types; a decoding error
 may occur before structured diagnostics can be generated.
 \item[Invalid model] Consult \texttt{isValid}, severity, code, message, and
 diagnostic path when included in the output. Correct the input and validate again.
 \item[No feasible solution] The structure may be valid even though constraints
 or demand cannot be satisfied. Review bounds and balances.
 \item[Unavailable backend] Check which families support it and its
 configuration. Validate only does not provide an alternative solver result.
 \item[Unexpected result] Review units, objective sign, state order,
 constraints, options, and exactness. Compare with simple calculations before
 changing algorithms.
\end{description}
The CLI exits successfully when it completes a command and unsuccessfully when
an error is thrown, including incorrect usage. Some validations return a report
with \texttt{isValid: false} normally; the exit code alone does not prove that
the model is valid. An automated workflow must also inspect the validation
document and result status. The absence of syntax errors does not demonstrate
that the model accurately represents the real system.
